



\section{K-theory of Stefanich rings}
\label{sec:stefanich}

\subsection{Successive K-theories}


For $r\geq0$, the endomorphisms of the monoidal unit
define a functor
\[
 \Omega_r:
 \operatorname{CAlg}(\nPr{(r+1)})\longrightarrow\operatorname{CAlg}(\nPr r),
 \qquad
 C\longmapsto\PrHom_C(\one_C,\one_C).
\]
The commutative algebra structure is induced by the
symmetric monoidal structure of $C$.

\begin{proposition}\label{prop:k-unit-end}
For $n\geq1$ and
$C\in\operatorname{CAlg}(\nPr{(n+1)})$, there is a natural
equivalence of commutative algebras
\begin{equation}\label{eq:k-unit-end}
 \gamma_{n,C}:
 \Omega_{n-1}K^{[n+1]}(C)
 \xrightarrow{\sim}
 K^{[n]}(\Omega_n C).
\end{equation}
\end{proposition}

\begin{proof}
The unit $u=\one_C$ is compact, since the
algebra unit $\one_{n+1}\to C$ preserves
compact objects.
By Corollary~\ref{cor:K-normalization},
the monoidal unit of $K^{[n+1]}(C)$
corresponds to $U_Cu$.
Theorem~\ref{thm:mapping-comparison} therefore
gives
\[
 \begin{aligned}
 \Omega_{n-1}K^{[n+1]}(C)
 &\simeq
 \PrHom_{\M_n(C)}(U_Cu,U_Cu)\\
 &\simeq K^{[n]}(\PrHom_C(u,u))
 =K^{[n]}(\Omega_n C).
 \end{aligned}
\]
For $n=1$, the middle equivalence is the
ordinary $K$-theory mapping formula.
Compatibility with commutative algebra
structures follows from the multiplicative
mapping comparison.
\end{proof}


A \emph{Stefanich ring} in our stable coefficient
tower is a coherent system
\[
 \mathbf R=(R_r,\alpha_r)_{r\geq0},
 \qquad
 R_r\in\operatorname{CAlg}(\nPr r),\qquad
 \alpha_r:R_r\xrightarrow{\sim}\Omega_rR_{r+1}.
\]
Morphisms are coherent families of commutative
algebra maps compatible with the $\alpha_r$.
Write $\operatorname{StRing}$ for the resulting
inverse limit category.
This is the stable, compactly generated variant
of \citep[Definition~3.13]{ScholzeGeometry}.


\begin{theorem}\label{thm:stefanich-k}
Higher K-theory induces a functor
\[
 K_{\mathrm{St}}:
 \operatorname{StRing}\longrightarrow
 \operatorname{StRing},
 \qquad
 K_{\mathrm{St}}(\mathbf R)_r
 =K^{[r+1]}(R_{r+1}).
\]
The structure equivalences are induced by
\eqref{eq:k-unit-end}.
\end{theorem}

\begin{proof}
Corollary~\ref{cor:K-normalization} gives
$K^{[r+1]}(R_{r+1})\in\operatorname{CAlg}(\nPr r)$, and
Proposition~\ref{prop:k-unit-end} gives
coherent equivalences
\[
 \begin{aligned}
 K^{[r+1]}(R_{r+1})
 &\xrightarrow{K^{[r+1]}(\alpha_{r+1})}
 K^{[r+1]}(\Omega_{r+1}R_{r+2})\\
 &\xrightarrow{\gamma_{r+1,R_{r+2}}^{-1}}
 \Omega_r K^{[r+2]}(R_{r+2}).
 \end{aligned}
\]
Naturality yields the asserted functor
on coherent towers.
\end{proof}

In particular, the zeroth level is
\(K_{\mathrm{St}}(\mathbf R)_0
 \simeq K(R_1^\omega)\),
which need not agree with $K(R_0)$.

\subsection{Examples}

By \citep[Lemma~2.11]{Aoki}
(see also \citep[Theorem~3.11]{ScholzeGeometry}),
the functor
\[
 \Mod_{(-)}(\nPr r):
 \operatorname{CAlg}(\nPr r)\longrightarrow
 \operatorname{CAlg}(\nPr{(r+1)})
\]
is fully faithful, with right adjoint $\Omega_r$.
In particular, \(\Omega_r\Mod_A(\nPr r)\simeq A\).




\begin{example}\label{ex:module-stefanich-k}
Let $a$ be a commutative ring spectrum. Define
\[
 R_0(a)=a,\qquad
 R_{r+1}(a)=\Mod_{R_r(a)}(\nPr r).
\]
These form a Stefanich ring, and
\[
 \begin{aligned}
 K_{\mathrm{St}}(\mathbf R(a))_0
   &\simeq K(a),\\
 K_{\mathrm{St}}^2(\mathbf R(a))_0
   &\simeq
   K_{\mathrm{MGS}}^{(2)}
     \bigl(\mathcal C_{R_2(a)}\bigr),
 \end{aligned}
\]
where $\mathcal C_{R_2(a)}$ is the small stable
$2$-category associated with $R_2(a)$ in
Corollary~\ref{cor:general-MGS}.
The second equivalence follows from that corollary
and Corollary~\ref{cor:K-normalization}.

For $a=\mathbb S$, the zeroth level is nonzero,
and hence so are all subsequent levels.
\end{example}


\begin{example}\label{ex:nonmonogenic-stefanich-zero}


Let $k$ be a field of characteristic zero,
viewed as an $\mathbb E_\infty$-ring, and put
\[
 R_0=k,\qquad
 R_1=\Fun(BC_2,\Mod_k),
\]
with pointwise tensor product, where $BC_2$
denotes the classifying space of the cyclic
group of order two.
Then $\Omega_0R_1\simeq k$.
Extend the tower by
$R_{r+1}=\Mod_{R_r}(\nPr r)$ for $r\geq1$.
Since $k[C_2]\simeq k\times k$, we obtain
\[
 K_{\mathrm{St}}(\mathbf R)_0
 \simeq K(k)\oplus K(k),
 \qquad
 K(R_0)\simeq K(k).
\]
Thus the two spectra are not equivalent,
as their $\pi_0$ have different ranks.
\end{example}



\begin{example}
Let $\mathbf R$ be a Stefanich ring.
For every $r\geq0$, the structure equivalence
\[
 K^{[r+1]}(R_{r+1})
 \xrightarrow[\sim]{
   \gamma_{r+1,R_{r+2}}^{-1}
   \circ K^{[r+1]}(\alpha_{r+1})}
 \Omega_rK^{[r+2]}(R_{r+2})
\]
induces, by adjunction, a natural morphism
\[
 \Mod_{K^{[r+1]}(R_{r+1})}(\nPr r)
 \longrightarrow K^{[r+2]}(R_{r+2}).
\]
In particular, the structure equivalence
$\alpha_0:R_0\simeq\Omega_0R_1$ induces
$\Mod_{R_0}(\Sp)\to R_1$, and hence
\[
 K(R_0)\longrightarrow K(R_1^\omega)
 \simeq(K_{\mathrm{St}}(\mathbf R))_0.
\]
Iterating these comparisons gives natural maps
\[
 K^m(R_0)\longrightarrow
 \bigl(K_{\mathrm{St}}^m(\mathbf R)\bigr)_0,
 \qquad m\geq0,
\]
where the superscripts denote iteration.
These maps need not be equivalences.
\end{example}







\subsection{Localizing motives and further directions}

The constructions of this paper suggest a parallel
theory of higher localizing motives. One approach
is to define them recursively, using localizing
motives at the preceding coefficient level.
By Lemma~\ref{lem:outgoing-coinverter},
the resulting
motives admit a presentation as localizations of
$\mathsf P_n(A)$. 
The main remaining issue is to
identify a satisfactory intrinsic notion of higher
Verdier exactness, compatible with the recursive
coefficient structure.

In a recursive Hom-wise model, a localizing analogue of
Theorem~\ref{thm:mapping-comparison} follows formally.
Nisnevich descent is also available. 
These constructions also provide a framework for
higher nonconnective K-theory and its analogues
for coherent Stefanich towers.

Another direction is the construction of higher
traces and Chern characters. In particular, one
may seek higher analogues of Hochschild homology,
topological Hochschild homology, and topological
cyclic homology, together with natural trace and
cyclotomic trace maps from higher K-theory.
A central question is whether these constructions
can retain the relevant noninvertible higher
morphisms and satisfy the required additive or
localizing properties. Such refinements could
provide new computational tools for higher
K-theory and its applications.
